\section{Sua parte: resultados, diagnóstico e próximos passos}
\textbf{Mensagem principal:} o modelo explica \nRdoisA{} da variação; renda em U, \% preta ou parda, urbanização e desigualdade se associam à violência; a variância não constante é a principal pendência.

\subsection{O que enfatizar}
\begin{itemize}
\item \textbf{Efeitos em \%:} como a resposta está em log, cada efeito é a variação da taxa para a mudança indicada, com as outras variáveis constantes. As cinzas não são significativas.
\item \textbf{Curva da renda:} pico em R\$ \nPico{} quando as outras variáveis estão controladas (diferente do pico dos dados brutos, que não controla nada).
\item \textbf{Diagnóstico:} normalidade razoável ($W$ = \nSW), \nNout{} outliers mantidos, vizinhos independentes (Durbin-Watson = \nDW), variância não constante (Breusch-Pagan = \nBP).
\item \textbf{Próximos passos:} seleção de variáveis (cap. 8), decidir as frágeis, tratar a variância, testar frequência escolar e garimpo.
\end{itemize}

\subsection{Perguntas prováveis}
\begin{description}[leftmargin=0pt]
\item[Por que \% intenção indeterminada tem efeito negativo?] Ela mede subnotificação: onde muitas mortes ficam como indeterminadas, parte dos homicídios some do registro. É um controle de qualidade do registro, não um fator de risco.
\item[A variância não constante invalida o modelo?] Não invalida os coeficientes (continuam certos na média), mas deixa os valores p aproximados; está nos próximos passos.
\item[Fronteira não ser significativa contradiz a hipótese do tráfico?] Com o estado no modelo, boa parte do efeito regional já é capturada pelo estado; a seleção de variáveis vai decidir se a fronteira fica.
\end{description}

\subsection{Encerramento}
Agradecer e abrir para perguntas. A tabela completa de coeficientes está no apêndice, depois do ``Obrigado''.
